\section{Differential equations as a first-class workload}
\label{sec:ode}

The differential-equation stack is where the three headline mechanisms
compound: pluggable oracles make input models declarative
(\secref{sec:ode-inputs}); protocols make solver families interchangeable
positions (\secref{sec:ode-solvers}); and structure preservation keeps
discretized, lifted, and iterated solver skeletons small and serializable
(\secref{sec:pde}).  \figref{fig:ode-taxonomy} maps the territory;
Table~\ref{tab:ode-coverage} compares coverage against existing frameworks.

\subsection{Input models: the problem layer stays native}
\label{sec:ode-inputs}

Every ODE method consumes one of four \emph{problem objects}, each holding
its generator as open-oracle-capable objects rather than matrices:

\begin{itemize}
\item \texttt{QODEProblem}: $u' = G u$ with the generator $G$ given as a
  block-encoding, the initial state as a state-preparation oracle, an
  optional \emph{dissipativity} promise with evidence provenance, and the
  initial norm;
\item \texttt{LinearODE}: $u' = -Au$ with Hermitian parts $A = L + iH$
  (the decomposition some methods need, declared by the caller);
\item \texttt{PolynomialODE}: $u' = \sum_p F_p\, u^{\otimes p}$ with each
  $F_p$ a block-encoding---the \emph{nonlinear} input model, consumed by
  Carleman lifting;
\item \texttt{FokkerPlanckProblem}: one-dimensional It\^o density dynamics
  $p' = -\partial_x(ap) + \partial_x^2(Dp)$, discretized by zero-flux finite
  volumes into a generator block-encoding, with classical reference
  witnesses (evolution, stationary and Boltzmann distributions) attached.
\end{itemize}

This is the problem--adapter--protocol discipline of
\secref{sec:requirements} made concrete: a PDE or ODE arrives in the
paradigm native to the application; adapters construct block-encodings
explicitly; solvers are generator protocols.  The same problem object feeds
every solver family, which is precisely what makes method comparison a
one-line change (\lstref{lst:qode-swap}).

\begin{lstlisting}[style=pyqpython, caption={One problem, two solver families (verbatim from the doctested tutorial).  Both generators return one-qubit state oracles for $u'=-u$, $t=0.05$; their ancilla structure, approximation route, and recovery conditions differ.  The tutorial is explicit that this checks \emph{assembly}: a degree-one Taylor kernel does not certify an arbitrary-precision solution---promise provenance (R6) records that boundary.}, label={lst:qode-swap}]
from pyqecclang import Bits, Builder, QODEProblem, identity, scale
from pyqecclang.algorithms.ode import linear_qode
from pyqecclang.algorithms.hamiltonian import taylor_hamiltonian
from functools import partial

b = Builder("initial_one", {"q": Bits(1)})
b.x(b["q"])
problem = QODEProblem(scale(-1, identity(1)), b.finish(),
                      dissipative=True, initial_norm=1.0)
kernel = partial(taylor_hamiltonian, degree=1)

lchs = linear_qode("lchs", hamiltonian_function=kernel)
schrodinger = linear_qode("schrodingerization", hamiltonian_function=kernel)

assert lchs.check(problem, time=0.05).ok
assert schrodinger.check(problem, time=0.05).ok
first = lchs.solve(problem, 0.05)
second = schrodinger.solve(problem, 0.05)
assert first.width == second.width == 1
\end{lstlisting}

\subsection{Six solver families as interchangeable protocols}
\label{sec:ode-solvers}

\paragraph{LCHS (linear combination of Hamiltonian simulations).}
Dissipative evolution $e^{Gt}$ is expanded as a quadrature sum of
Hermitian-branch evolutions $e^{-i(H + kL)t}$ weighted by Cauchy-type
coefficients; the library assembles the weighted LCU~\cite{an2026lchs} with a configurable
\texttt{QuadraturePlan} (default: finite Cauchy cutoff and spacing).  The
finite-quadrature remainder is recorded as a pending item on the generated
modules---an R6 obligation, not a silent truncation.

\paragraph{Schr\"odingerization.}
For a general generator $G = H_1 + iH_2$, the state is lifted through a
Fourier-transformed ancilla register so that evolution becomes Hamiltonian
in the lifted space; the library realizes the lift with a
Fourier-momentum diagonal block-encoding, warps the initial state by
$e^{-|p|}$, and selects the physical channel at a planned index with an
explicit recovery scale~\cite{jin2024schrodingerization,hu2024schrocircuits}.  Periodic-truncation caveats are recorded as
pending.

\paragraph{Contour-based matrix-function decomposition (CBMD).}
Evolution (or a general matrix function) is decomposed along a contour~\cite{wang2026cbmd} with
auxiliary poles; the assembly supports arbitrary nodes and residues, and the
generated modules record which auxiliary non-Hermitian evolutions and series
tails were omitted.

\paragraph{Carleman linearization.}
The nonlinear path: \texttt{PolynomialODE} inputs are lifted to a cutoff-$K$
block-triangular tensor system (the lift itself is a structured placement of
trilinear blocks with level controls), the initial state becomes a weighted
tensor power, any linear solver protocol then applies, and the output is
projected back to the level-one subspace~\cite{liu2021pnas,liu2023carleman}.  The truncation tail is a recorded
pending item.

\paragraph{Implicit-Euler history systems via QLSS.}
Time discretization first, quantum solving second: the whole solution
history of the implicit-Euler scheme is assembled as one block system (an
LCU of projector, shift, and generator blocks) and solved in one shot by any
QLSS protocol; the final-time subspace is selected on the quantum side.  In
this family, \emph{the linear-system solver is the injectable position}.

\paragraph{Truncated Taylor kernels.}
A closed polynomial candidate $\sum_k (tG)^k/k!$ with no structural
assumptions---also the default \emph{Hamiltonian-function kernel} that the
three evolution-based families above inject, with degree and configuration
swappable at call time.

\begin{figure}[t]
  \centering
  \begin{tikzpicture}[
      font=\small,
      prob/.style={draw=black!60, rounded corners=2pt, fill=black!8, align=center, inner sep=3.5pt, text width=5.5cm},
      fam/.style={draw=black!60, rounded corners=2pt, fill=black!3, align=center, inner sep=3pt, text width=4.7cm},
      inj/.style={draw=black!45, rounded corners=2pt, fill=white, dashed, align=center, inner sep=3pt},
      grp/.style={draw=black!35, rounded corners=4pt, dotted},
      arr/.style={-{Stealth[length=2mm]}, black!60},
      darr/.style={-{Stealth[length=2mm]}, black!45, dashed},
      lbl/.style={font=\scriptsize\itshape, black!55},
    ]
    % Problem layer
    \node[prob] (pdeA) at (-3.5, -0.35)
      {\textbf{PDEs via spatial discretization}\\[-2pt]\scriptsize heat equation, Fokker--Planck};
    \node[prob] (pdeB) at (3.5, -0.35)
      {\textbf{nonlinear polynomial PDEs}\\[-2pt]\scriptsize Burgers, KdV, reaction--diffusion, 2-D vector Burgers};
    \node[prob] (probA) at (-3.5, -2.3)
      {\textbf{linear problem}\\[-2pt]\scriptsize $u'=Gu$: block-encoding of $G$, initial state, promises};
    \node[prob] (probB) at (3.5, -2.3)
      {\textbf{polynomial problem}\\[-2pt]\scriptsize $u'=\sum_p F_p u^{\otimes p}$: block-encodings of $F_p$};
    \draw[arr] (pdeA) -- (probA);
    \draw[arr] (pdeB) -- (probB);

    % Carleman branch
    \node[fam] (carleman) at (3.5, -4.1)
      {\textbf{Carleman lift}\\[-2pt]\scriptsize cutoff $K$, tensor placement};
    \draw[arr] (probB) -- (carleman);

    % Solver-family container
    \node[fam] (lchs)  at (-3.5, -5.5)  {\textbf{LCHS}\\[-2pt]\scriptsize weighted Hermitian branches};
    \node[fam] (schro) at (-3.5, -6.6)  {{\footnotesize\textbf{Schr\"odingerization}}\\[-2pt]\scriptsize Fourier lift, recovery channel};
    \node[fam] (cbmd)  at (-3.5, -7.7)  {\textbf{contour (CBMD)}\\[-2pt]\scriptsize nodes, residues, omitted tails};
    \node[fam] (euler) at (-3.5, -8.8)  {\textbf{Euler $\rightarrow$ QLSS}\\[-2pt]\scriptsize one-shot history block system};
    \node[fam] (taylor) at (-3.5, -9.9) {\textbf{Taylor kernel}\\[-2pt]\scriptsize $\sum_k (tG)^k/k!$};
    \node[grp, minimum width=6.1cm, minimum height=6.35cm] (box) at (-3.5, -7.7) {};
    \node[lbl, anchor=south west] at (-6.5, -4.4) {linear-solver families};
    \draw[arr] (probA.south) -- (box.north);
    \draw[arr] (carleman.west) to[out=180, in=40] ([yshift=2.45cm]box.east);
    \node[lbl, anchor=south east] at (1.55, -4.35) {any linear solver};

    % Injectable kernels
    \node[inj] (hfun) at (3.0, -5.85)
      {\textbf{injectable Hamiltonian-}\\\textbf{function kernel}\\[-2pt]\scriptsize default: truncated Taylor;\\[-2pt]\scriptsize degree per call};
    \node[inj] (qlssk) at (3.0, -8.5)
      {\textbf{injectable QLSS kernel}\\[-2pt]\scriptsize Costa adiabatic walk /\\[-2pt]\scriptsize CKS Chebyshev};
    \draw[darr] (hfun.west) -- (lchs.east);
    \draw[darr] (hfun.south west) -- (schro.east);
    \draw[darr] (hfun.south) to[out=-90, in=15] (cbmd.east);
    \draw[darr] (qlssk.west) -- (euler.east);
  \end{tikzpicture}
  \caption{The ODE/PDE solver territory of \pqlang{}.  Linear problems feed
  four evolution- and system-oriented families (LCHS, Schr\"odingerization,
  contour decomposition, Taylor kernels) plus the QLSS-based history-system
  family; nonlinear polynomial problems lift through Carleman into the same
  linear-solver position.  Dashed boxes mark \emph{injectable protocols}:
  the kernel realizing $e^{-iHt}$ subroutines and the linear-system kernel
  are swappable arguments, as are the spatial discretizers connecting the
  PDE layer to the problem layer.}
  \label{fig:ode-taxonomy}
\end{figure}

\subsection{The PDE layer: discretization, homotopy analysis, finite volumes}
\label{sec:pde}

Three paths connect PDEs to the ODE stack, each exercising a different part
of the language:

\paragraph{Pluggable spatial discretization.}
\texttt{DiscretePDE} pairs a generator block-encoding with an initial state
after spatial discretization; \texttt{qpde\_solver} threads any
discretizer into any QODE protocol.  The built-in periodic heat equation,
for instance, assembles a structured four-point shift block-encoding (the
stencil shift primitive of \secref{sec:be-algebra}) and runs under LCHS or
Schr\"odingerization interchangeably.

\paragraph{QHAM: homotopy-analysis linearization of nonlinear PDEs.}
Symbolic polynomial PDEs---Burgers, KdV, reaction--diffusion ($u' = u -
u^3$), coupled systems, and two-dimensional vector Burgers---are expressed
in a \emph{PolynomialPDE} layer; a homotopy-analysis plan (order parameter
$\eta$, linearization order $m$) is computed \emph{lazily} as a row-rule
plan that serializes without materializing its tensor blocks; bindings then
generate the lifted generator block-encoding and the norm-preserving lifted
initial state; and the result is handed to any QODE solver, with
\texttt{dissipative\_shift} adapting it for LCHS/CBMD.  Laziness here is an
R5 statement: a second-order QHAM plan for a moderately sized Burgers grid
is a \emph{plan object} with a schema, not an exploded circuit.

\paragraph{Fokker--Planck from SDEs.}
Stochastic dynamics enter as It\^o density evolution, are discretized by
zero-flux finite volumes into Pauli-structured generators, and are solved as
\texttt{QODEProblem}s; the problem object carries pure-Python classical
witnesses (evolution, stationary and Boltzmann distributions) used by tests
to hold the input model honest.

\paragraph{A note on method coverage.}
We are not aware of another programming system in which these solver
families coexist behind a common problem-object interface with injectable
linear-system and Hamiltonian-function kernels; the comparison of
Table~\ref{tab:ode-coverage} makes this precise.  The closest analogues are
individual research codebases (one method each, no shared input model) and
textbook circuit libraries (linear systems as demonstrations, no ODE/PDE
stack at all).  Coverage, of course, is only meaningful together with
verification status, which \secref{sec:implementation} reports per family.
